\documentclass[10pt]{article}
\usepackage[margin=20mm]{geometry}
\usepackage{amsmath,amssymb,mathtools,fontspec,unicode-math,xurl,hyperref}
\IfFontExistsTF{DejaVu Serif}{\setmainfont{DejaVu Serif}}{\setmainfont{Latin Modern Roman}}
\setmathfont{latinmodern-math.otf}
\setlength{\parindent}{0pt}\setlength{\parskip}{6pt}\sloppy
\title{Birational descent, local triviality, and two auxiliary estimates}
\author{AI-assisted technical review}
\date{}
\begin{document}
\maketitle
Abstract: Lemma 4.5 is false under its stated hypotheses. A product over the blow-up of the projective plane, with the identity map on total spaces and the blow-down on bases, satisfies every hypothesis but has a dimension jump in the descended fibers. The proof also makes an invalid assertion about curves lying over a contracted base curve. The lemma is invoked directly in Theorem 4.9 and in the final proof of Theorem 1.1, so this is a Major defect in the written proof chain. A second Major proof-inference issue is the canonical-bundle-free vanishing used in Lemma 4.4: a smooth plane quartic gives a nonzero cohomology group under the stated choices. A corrected choice of the auxiliary ample divisor repairs that vanishing step. A missing constant in Lemma 2.8 is Minor. No counterexample to the principal structure theorem is claimed.\par
\section*{Exact source and review scope}
The bound source is Shin-ichi Matsumura and Juanyong Wang, Structure theorem for projective klt pairs with nef anti-canonical divisor, arXiv:2105.14308v3 [paper]. The upload timestamp is 30 August 2023, 03:32:48 UTC; the supplied 42-page PDF internally prints Date: August 31, 2023, version 0.01. Its SHA256 is recorded in full in provenance and inventory. The associated journal DOI is 10.4171/JEMS/1702. The publisher records submission on 25 January 2024, acceptance on 6 July 2025, and online publication on 30 September 2025 [journal]. The arXiv record does not identify v3 as the accepted journal text. The findings below apply to the exact supplied preprint; equivalence with the later journal article was not established.\par
All 42 pages were read as extracted text, including the references. The official version-specific HTML was also consulted and retained for the three disputed passages. Its generated date in 2026 is not treated as a new author revision. The source PDF remains governing. All 30 required source evidence pages were actually displayed and inspected in the same-chat continuation, including the birational and universal-cover diagrams, the exact assumptions of Lemma 4.5, the vanishing formula, and their uses in the main proof. The earlier `EMFILE` failures and false viewing records are preserved. This visual check does not alter the three mathematical findings.\par
\section*{The main logical route}
The theorem concerns complex projective klt pairs with effective rational boundary and nef \(- (K_X+\Delta)\). It allows a finite quasi-etale cover before obtaining a holomorphic, locally constant MRC fibration. Locally constant is defined through a monodromy representation on a fixed fiber over the universal cover, and is stronger than merely locally trivial. The theorem does not claim that the original variety already has this structure without a cover. The fundamental-group statements in Section 1.3 are explicitly conjectures and are not used as proved inputs.\par
The initial reduction takes a maximally quasi-etale cover and a crepant terminal model. Terminality supplies smoothness in codimension two, which is used in Theorem 3.8. For each normalized direct image, Proposition 3.9 converts a further quasi-etale cover into an etale one by maximality and descends local freeness. This handles the different sheaves without an unproved simultaneous infinite tower construction.\par
Theorem 3.2 descends the selected flat connection along rationally connected general fibers, then extends it across a general point of each relevant divisor using a reduced component and a local section. The normalization of the graph, rather than the whole resolution exceptional set, is important: only the former exceptional set is known not to dominate the base. Proposition 3.3(d) supplies compatibility of multiplication maps, which is required by Proposition 2.5 and cannot be dropped on a noncompact base. The subsequent foliation argument gives the birational structure. The obstruction identified below occurs in the later passage from that structure to a locally trivial fibration.\par
\section*{A false birational descent lemma}
Problem 1 (Major: Lemma 4.5, pages 29–30).  The assertion that \(f\) must be locally trivial does not follow from the hypotheses stated in Lemma 4.5. There is an explicit smooth projective counterexample.\par
\subsection*{Construction and complete hypothesis check}
Let \[W=\mathbb P^2,\qquad b:W'=\operatorname{Bl}_p\mathbb P^2\longrightarrow W,\qquad
 V'=V=W'\times\mathbb P^1.\] Put \[\beta_W=b,\quad \beta_V=\operatorname{id}_{V'},\quad
 f'=\operatorname{pr}_{W'},\quad f=b\circ f',\quad \Delta'=0 .\] The square commutes. All four varieties are smooth, normal, and projective. Both horizontal maps are projective birational morphisms; the identity is a birational morphism. Both vertical maps are proper surjective morphisms with connected fibers, exactly the meaning of fiber space used in the paper. No equidimensionality requirement is included in that definition or in Lemma 4.5.\par
The map \(f'\) is a product fibration with fiber \(\mathbb P^1\), hence is locally constant and has simply connected fiber. The base \(W\) and the pair \((V',0)\) are smooth and therefore klt. Finally, \[K_{V'/W'}=\operatorname{pr}_{\mathbb P^1}^{*}K_{\mathbb P^1},\qquad
 - (K_{V'/W'}+\Delta')=2\operatorname{pr}_{\mathbb P^1}^{*}\mathcal O_{\mathbb P^1}(1)\] is nef. This verifies the precise relative, rather than absolute, positivity hypothesis in the lemma.\par
If \(q\ne p\), then \(f^{-1}(q)\simeq\mathbb P^1\), of dimension one. At the blown-up point, \[f^{-1}(p)=b^{-1}(p)\times\mathbb P^1\simeq\mathbb P^1\times\mathbb P^1,\] of dimension two. A locally trivial analytic fibration cannot have that dimension jump. This is a counterexample to the lemma’s conclusion, not merely a concern about an omitted justification.\par
For a local algebra check, use a blow-up chart with \[(u,v,w)\longmapsto(x,y)=(u,uv),\] where \(w\) is an affine coordinate in the second \(\mathbb P^1\). The fiber ideal at \((0,0)\) is \((u,uv)=(u)\), leaving two free coordinates. At \((1,0)\) it is \((u-1,uv)=(u-1,v)\), leaving one. These reductions and the Jacobian ranks were independently executed with exact symbolic arithmetic. The global product construction supplies properness and the second blow-up chart; these are not inferred solely from the local computation.\par
\subsection*{The precise failed inference}
On page 30, in the paragraph beginning with an ample divisor on \(V\), the proof chooses a curve \(C_1\) above a curve \(C_0\) contracted by \(\beta_W\), and asserts that \(C_1\) must be contracted by \(\beta_V\). Containment of \(\beta_V(C_1)\) in a fiber of \(f\) does not imply that \(C_1\) is contracted by \(f'\).\par
In the example, let \(E=b^{-1}(p)\) and \(C_1=E\times\{z_0\}\). Then \(f'(C_1)=E\), while \(\beta_V\) is the identity and does not contract \(C_1\). The claimed zero intersection can also be contradicted with an ample divisor. With \(H=b^*\mathcal O_{\mathbb P^2}(1)\), the divisor \(2H-E\) is ample on \(W'\); its intersections are \[(2H-E)\cdot E=1,\qquad (2H-E)^2=3.\] Thus \(A=\operatorname{pr}_{W'}^*(2H-E)+
\operatorname{pr}_{\mathbb P^1}^*\mathcal O_{\mathbb P^1}(1)\) is ample on \(V\), yet \(\beta_V^*A\cdot C_1=1\), not zero.\par
\subsection*{Dependence and repair limits}
Lemma 4.5 is invoked on page 36 in Step 3 of Theorem 4.9 and on page 37 in Step 2 of Theorem 1.1. Its failure therefore affects a core descent step. The counterexample does not disprove Theorem 1.1: the theorem asserts existence after a cover, and the counterexample is directed at an auxiliary statement about a specified diagram.\par
The applications have additional geometric information, including equidimensionality established earlier. It may permit a narrower replacement lemma. This report does not claim that adding equidimensionality alone proves the desired descent, nor that the original curve argument becomes correct after adding it. A repaired argument must justify the necessary contraction or descent property in the actual diagrams, check the pair and boundary, and then revisit both uses of Lemma 4.5. The invalid general lemma cannot simply be cited.\par
\section*{Vanishing without the canonical twist}
Problem 2 (Major proof inference: Lemma 4.4, pages 28–29).  The final part of the proof asserts \[H^i\bigl(W,E_m\otimes\mathcal O_W((d+1-i)A)\bigr)=0\quad(i>0)\] for a very ample \(A\), by Nakano vanishing. Here \(d\) must mean the dimension of \(W\); the source writes \(Y\). The asserted vanishing without a canonical twist or a stronger choice of \(A\) is false, even for trivial \(E_m\).\par
Let \(C\subset\mathbb P^2\) be a smooth plane quartic, for example the Fermat quartic over \(\mathbb C\). By adjunction, \[g(C)=3,\qquad K_C\simeq\mathcal O_C(1).\] Take \[W=C,\quad V=C\times\mathbb P^1,\quad h=\operatorname{pr}_C,\quad D=0,\quad
 B=h^*\mathcal O_C(1)+\operatorname{pr}_{\mathbb P^1}^*\mathcal O_{\mathbb P^1}(1).\] This is a locally constant product with irregularity-zero fiber; \((V,D)\) is klt, the relative anti-canonical divisor is nef, and \(B\) is ample, hence nef and big. The normalized decomposition in Lemma 2.4 may be taken as \(B_0=\operatorname{pr}_{\mathbb P^1}^*\mathcal O(1)\), \(B_W=\mathcal O_C(1)\). Consequently \[E_m=h_*\mathcal O_V(mB_0)\simeq\mathcal O_C^{\oplus(m+1)}.\] Choose the very ample divisor \(A=\mathcal O_C(1)=K_C\), which satisfies the choice made in the source. For \(d=i=1\), Serre duality gives \[h^1(C,E_m\otimes A)=(m+1)h^1(C,K_C)=m+1>0 .\] Thus all of the lemma’s geometric hypotheses and the displayed positivity of the twisted trivial bundle hold, but the cited vanishing does not.\par
This inference supplies Castelnuovo–Mumford regularity, uniform global generation, and ultimately nefness of \(B_0\). Lemma 4.4 is then used in Lemma 4.5 and explicitly in Theorem 4.12 on page 36. The Major classification concerns that written core inference; it does not assert that Lemma 4.4’s conclusion is false.\par
A local repair of this particular step is available. Choose a single very ample \(A\) with \(A-K_W\) ample, after the reduction to smooth \(W\). For \(1\le i\le d\), \((d+1-i)A-K_W\) is ample. Use the already established numerical flatness of \(E_m\), not an unqualified assertion about arbitrary flat bundles: its filtration by unitary flat quotients permits Kodaira–Nakano vanishing for each quotient twisted by this ample line bundle, and the resulting long exact sequences give the required vanishing for \(E_m\). This uses the same fixed \(A\) for every \(m\).\par
The relative basepoint-free step on page 29 should also use the klt pair \((V,D)\), the nefness of \(- (K_{V/W}+D)\), and sufficiently divisible multiples of the relatively nef and big \(B_0\). The hypothesis does not assert that \(-K_{V/W}\) itself is nef. This supplies the intended type of closing argument under the earlier numerical-flatness conclusions; it does not repair the false Lemma 4.5.\par
\section*{A boundedness estimate with a missing factor}
Problem 3 (Minor: Lemma 2.8, page 14).  The displayed estimate for \(e=\sum_{i=1}^m s_i\otimes f_i\) omits a factor \(m\).\par
For the identity morphism of a smooth disk, the trivial line bundle with its trivial metric, and \(s_1=s_2=f_1=f_2=1\), every hypothesis of Lemma 2.8 holds. The written inequality becomes \(4\le2\). The valid estimate is \[\left|\sum_i f_i s_i\right|^2
       \le m\sum_i |f_i|^2|s_i|^2.\] Alternatively one can use the product of the two sums of squared norms. This fixed local factor does not change local upper boundedness or the plurisubharmonic extension argument. It is therefore Minor and does not invalidate Lemma 2.8’s conclusion. Local generation of the pullback can be implemented after shrinking the neighborhoods; no global section-tensor identity on arbitrary large analytic open sets is needed.\par
\section*{Further checks and limits}
The quasi-etale reductions were checked at the stated level: normality and codimension-two complements enter the fundamental-group diagram in Theorem 2.12; the selected terminal model remains maximal. The example in Section 3.3 correctly warns that a quasi-etale cover can change the MRC fiber dimension. It is therefore essential to recompute the fibration after taking the cover, as the construction does.\par
The singular-metric discussion distinguishes locally free pullback from torsion-free pullback and supplies reflexivity and codimension hypotheses in Lemma 2.8. The metric comparison in Lemma 2.10 requires its exceptional potential adjustment; a curvature lower bound with a resolution metric cannot be pushed down unchanged as a bound by a fixed metric on the base. The direct-image normalization uses \(m/r\in\mathbb Z\) and a fixed large exceptional coefficient. Those restrictions were retained in checking Theorem 3.2.\par
The connection descent uses smooth rationally connected general fibers, so the needed vanishing of one-forms is applied there. No defect is asserted from a possible convention concerning torsion in differential sheaves on singular varieties in the introductory Theorem 2.2. Likewise, the selected connection associated with the numerical-flatness construction is not identified with every possible flat connection on the same holomorphic bundle.\par
In the local-constancy criterion, compatibility of the graded multiplication maps and the boundary ideal is substantive. Remark 2.6(c) reduces the needed data to finitely many degrees; the proof does not need removal of an uncontrolled countable union of codimension-two sets. The relative anti-canonical theorem and the later foliation results depend on several external positivity, semistable reduction, and integrability theorems. Their use and advertised hypotheses were traced, but those external proofs and the modifications described as carrying over are not independently certified here. The Calabi–Yau product consequence on page 4 likewise depends on the cited Ambro and Druel inputs.\par
The full coverage ledger records pages 1–7 for framing, 7–17 for preliminaries, 17–25 for flat direct images, 25–31 for local constancy and birational descent, 31–37 for the main proof, and 37–42 for references. A bounded search located the exact arXiv record, the publisher record, the official v3 HTML, and related earlier work. No verified separate supplement, paper-specific code, or erratum was retrieved. No author code was executed.\par
The overall supported assessment is Major. The exact computations are supporting checks for the analytic and algebraic proofs above, not a substitute for them. The previous report and incomplete delivery records are retained. This revised TeX requires fresh compilation and actual inspection of every resulting PDF page. Completion is bound to the exact current source, TeX, compiler, retained dependencies, PDF and observed image hashes in the continuation provenance. The scientific decision is `pending` and human-review fields remain empty; this is an initial review, not independent acceptance.\par
9 S.-i. Matsumura and J. Wang, Structure theorem for projective klt pairs with nef anti-canonical divisor, arXiv:2105.14308v3, 2023, https://arxiv.org/abs/2105.14308v3. Official version-specific HTML, https://arxiv.org/html/2105.14308v3, secondary corroboration retained with retrieval metadata. EMS Press, DOI 10.4171/JEMS/1702, article record, https://ems.press/journals/jems/articles/14299169.\par
\end{document}
